\subsection{Plan de Aceptación del Producto}

Los criterios de aceptación de cada entregable son acordados entre el equipo y el director desde el inicio del proyecto. Todos los criterios deben ser verificables objetivamente; ninguno depende de apreciaciones subjetivas del evaluador. Cada uno nombra el comando, el archivo o el enlace con que se verifica, y lo que no se cumplió se dice.

\subsubsection{Prototipo funcional del SDK}

El SDK se considera aceptado cuando cumple cuatro condiciones. Compila sin errores en TypeScript estricto (\texttt{"strict": true} en \texttt{sdk/tsconfig.json}). Las pruebas de los cuatro adaptadores (Wompi, Rapyd, Mercado Pago y Kushki) pasan en la integración continua, y sus pruebas de contrato pasan contra los sandboxes reales (\texttt{npm run test:sandbox}). La cobertura de pruebas es igual o superior al 80\% en ramas, funciones, líneas y sentencias, medida sobre todo el paquete y no solo sobre el núcleo (\texttt{coverageThreshold} de \texttt{sdk/jest.config.js}). Por último, el paquete publicado se instala desde npm y se carga en un proyecto nuevo, en CommonJS y en ES Modules, sin modificar su código fuente. Lo comprueba \texttt{npm run check:published}, y la integración continua corre \texttt{npm pack --dry-run} en cada cambio. Las herramientas de validación son Jest, \texttt{npm pack} y \texttt{check:published}.

\subsubsection{API de simulación desplegada}

La API se considera aceptada cuando cumple cuatro condiciones. Reproduce, para cada pasarela, los desenlaces de aprobación, rechazo, pendiente, tiempo de espera y error, y emite los webhooks firmados, todo alcanzable desde el SDK sin cabeceras especiales. Su cobertura de pruebas es igual o superior al 80\% con el mismo criterio del SDK (\texttt{simulator-api/jest.config.js}). Está desplegada en Render y es accesible desde internet sin credenciales reales (\url{https://kit-pagos-colombia.onrender.com}). Por último, la colección de Postman versionada en \texttt{simulator-api/postman/} pasa al ejecutarse con Newman contra el simulador local y contra Render (\path{docs/02-arquitectura/4-despliegue-en-render-y-coleccion-postman.md}, sección 6). La versión 1.0 de este plan pedía además una respuesta en menos de 500~ms; ese tiempo no se midió, y el SRS v2.0 también lo quitó. Las herramientas de validación son Jest, Postman con Newman y los registros del servidor de despliegue.

\subsubsection{Documentación centralizada de datos de prueba}

La documentación se considera aceptada cuando cubre las cuatro pasarelas. Por proveedor debe traer al menos tarjetas de prueba con sus CVV y fechas de vencimiento, documentos de identidad válidos para sandbox, bancos PSE y referencias de pago de ejemplo. Debe estar publicada en el repositorio público, en \texttt{docs/testing-data/}. Cada afirmación declara su nivel de evidencia: medida contra el sandbox con fecha, medida solo en parte del flujo o tomada de la documentación oficial sin comprobar (\texttt{docs/testing-data/README.md}, sección 3). Lo medido lo respaldan las pruebas de contrato contra los sandboxes.

\subsubsection{Informe comparativo de la evaluación}

El informe se considera aceptado cuando contiene las seis variables del experimento, medidas sobre el prototipo B, que integra las pasarelas con el framework, y el prototipo D, que lo hace directamente contra las APIs nativas. Entre ellas están las métricas WMC, CBO y RFC, calculadas con ts-morph y adaptadas a archivos porque el código de pagos de B no tiene clases (\path{docs/project-management/prototypes-evaluation-plan.md}; \texttt{methodology.md}, sección 6.1.1). Los prototipos deben estar en \texttt{prototypes/} con el detalle suficiente para que el experimento sea reproducible, y el análisis comparativo debe interpretar las diferencias observadas en términos de complejidad y acoplamiento. Las mediciones están en \texttt{docs/evaluation/} desde el 28 de octubre de 2026, y el informe comparativo, desde el 9 de noviembre, cuando se alcanzó el hito H5; falta el capítulo VII de la memoria. El director valida la pertinencia del análisis.

\subsubsection{Documentación final del trabajo de grado y repositorio público}

La documentación final se considera aceptada cuando cumple cuatro condiciones. Cubre todas las fases del ciclo DSR con sus resultados. Están todos los documentos requeridos (la memoria, el SRS, el SAD, el plan de pruebas y este SPMP), en LaTeX en \texttt{docs/final-doc/} y compilados por la integración continua (\path{.github/workflows/final-doc.yml}). El repositorio público incluye instrucciones de instalación (\texttt{sdk/README.md}), ejemplos de integración para cada pasarela (\texttt{examples/}) y una página de presentación publicada en GitHub Pages (\path{.github/workflows/deploy-landing.yml}). Por último, un integrante que no escribió una sección puede seguir las instrucciones del repositorio e integrar el SDK en un proyecto nuevo sin contactar a los autores. Esta prueba no se hizo, así que ese criterio queda sin cumplir. La versión 1.0 de este plan pedía también un SDD y documentación generada con TypeDoc; ninguno de los dos se produjo. El director aprueba formalmente el documento.
